\chapter[System Design, Data and Preprocessing]{System Design, Data\\ and Preprocessing}\label{ch:system}

\section{Architecture of the framework}
The framework is organised in layers that mirror the workflow in the proposal (\cref{fig:architecture}). Each
layer is a Python module with a single responsibility. The modules communicate only through files with stable
formats (Parquet, SQLite and JSON), so any layer can be replaced without changing the others. For example, the
Streamlit dashboard could be swapped for Power BI reading the same SQLite store.

\begin{figure}[htbp]\centering
\begin{tikzpicture}[node distance=3.5mm, every node/.style={font=\footnotesize},
  L/.style={draw=brand, rounded corners, fill=brand!7, minimum width=12.5cm, minimum height=8.5mm, align=left, text width=12cm},
  arr/.style={-{Stealth[length=2mm]}, thick, brand}]
\node[L] (a) {\textbf{1. Data sources} \hfill raw CSV export (\NRows transactions, \NCols columns)};
\node[L, below=of a] (b) {\textbf{2. Ingestion \& validation} \hfill typed load, schema, rule and identity checks \hfill \texttt{data.py}};
\node[L, below=of b] (c) {\textbf{3. Cleaning \& storage} \hfill imputation, de-duplication, outlier flags, star schema (SQLite) \hfill \texttt{data.py}};
\node[L, below=of c] (d) {\textbf{4. Exploratory analytics} \hfill trends, seasonality, tests, elasticity \hfill \texttt{eda.py}};
\node[L, below=of d] (e) {\textbf{5. Feature engineering} \hfill leakage-safe lags, rolling statistics, calendar, price \hfill \texttt{features.py}};
\node[L, below=of e] (f) {\textbf{6. Machine learning} \hfill forecasting $\cdot$ classification $\cdot$ clustering $\cdot$ anomaly detection};
\node[L, below=of f] (g) {\textbf{7. Evaluation \& explanation} \hfill benchmarks, metrics, permutation importance, SHAP};
\node[L, below=of g] (h) {\textbf{8. Presentation} \hfill Streamlit dashboard $\cdot$ project website $\cdot$ \LaTeX{} report \hfill \texttt{dashboard/}};
\node[L, below=of h, fill=accent!10, draw=accent] (i) {\textbf{9. Decision support} \hfill automatically generated insights and recommendations \hfill \texttt{insights.py}};
\foreach \x/\y in {a/b,b/c,c/d,d/e,e/f,f/g,g/h,h/i} \draw[arr] (\x)--(\y);
\end{tikzpicture}
\caption{Layered architecture of the intelligent analytics framework}\label{fig:architecture}
\end{figure}

A single command, \texttt{python run\_pipeline.py}, executes all layers in order. It writes:
\begin{itemize}
  \item the cleaned dataset (\path{data/processed}) and the warehouse (\path{data/warehouse/dancake.db});
  \item metrics in JSON (\path{outputs/metrics}) and tables in CSV (\path{outputs/tables});
  \item every figure and \LaTeX{} table of this report (\path{report/figures}, \path{report/tables}), and a
        file of \LaTeX{} macros holding every number quoted in the text (\path{report/generated/numbers.tex});
  \item compact aggregated data for the dashboard and the website (\path{data/dashboard}).
\end{itemize}
Because no result is typed by hand, the report, dashboard and website cannot drift apart. A new month of data
can be incorporated by re-running one command.

\section{The dataset}
\Cref{tab:dqsummary} summarises the dataset and the results of the validation layer. Each transaction belongs to
one of \NProducts products in 13 categories. These are sold across \NRegions regions through \NChannels
channels (General Trade, Modern Trade, Distributor, E-commerce) to \NCustTypes customer types (Retailer, Super
Shop, Distributor, Corporate, Online Customer). Five products were launched after 2020: Lachcha Semai, two
breads, a croissant and a cappuccino muffin gift pack.

\begin{table}[htbp]\centering\small
\caption{Dataset profile and data-quality validation results}\label{tab:dqsummary}
\begin{tabular}{@{}lr@{}}
\toprule
\textbf{Check} & \textbf{Result} \\ \midrule
Rows $\times$ columns & \NRows $\times$ \NCols \\
Period covered & \DateMin -- \DateMax \\
Products / regions / channels / customer types & \NProducts / \NRegions / \NChannels / \NCustTypes \\
Missing cells (incl.\ hidden placeholders such as ``NA'', ``-'') & 0 \\
Exact duplicate rows / duplicate transaction IDs & 0 / 0 \\
Domain-rule violations (quantity, price, discount, value) & 0 \\
Gross $=$ quantity $\times$ market price (max.\ abs.\ error) & $<10^{-11}$ BDT \\
Net $=$ gross $-$ discount (max.\ abs.\ error) & $<10^{-11}$ BDT \\
Year/month columns consistent with date & 100\% \\
Product ID $\rightarrow$ one name, category, pack size & 100\% \\
Statistical quantity outliers (flagged, retained) & \NOutlierFlags (\OutlierPct\%) \\
Warehouse reconciliation (SQL total $-$ pandas total) & 0.00 BDT \\ \bottomrule
\end{tabular}
\end{table}

\section{Preprocessing pipeline}
The pipeline is written defensively. It must work correctly on dirty data, not only on the delivered file,
which proved to be exceptionally clean. Its behaviour on dirty data was verified in two ways. Six hands-on
practice scripts deliberately damaged copies of the data and measured how much was repaired. In addition, the
automated tests in \texttt{tests/} inject missing values, duplicates and impossible values and assert that the
pipeline repairs them (\cref{app:reproduce}). The steps are applied in a fixed order, because some steps depend
on others.

\subsection{Data type conversion and validation}
Dates are parsed with an explicit format (\texttt{dd-mm-yy}) and numeric columns are coerced. Values that fail
to parse become missing, so they are never silently misread. Text fields are trimmed of spaces. The validation
step then produces a machine-readable quality report (\texttt{outputs/metrics/data\_quality.json}). The report
checks for missing values, hidden placeholder strings, duplicate rows and keys, and domain-rule violations. It
also checks the accounting identities between columns and the consistency of calendar fields and product
attributes.

\subsection{Missing-value treatment}
A decision hierarchy was implemented, following \citet{garcia2015}:
\begin{enumerate}
  \item \emph{derive} the value exactly from other columns where an identity exists (for example,
        $\text{quantity} = \text{gross}/\text{market price}$, or year and month from the date);
  \item \emph{drop} the row if an essential identifier (transaction ID, date, product) is missing, because a
        guess would amount to fabrication;
  \item \emph{impute} numeric gaps with the median of the same product, which is robust to skew;
  \item \emph{impute} product attributes with the product's mode, and label missing market attributes
        explicitly as ``Unknown'' so that they remain visible.
\end{enumerate}
In practice tests with 3\% of cells removed at random, deriving values from accounting identities restored
quantities, prices and sales values exactly. Simple mean imputation, by contrast, introduced errors of several
units per order.

\subsection{Duplicate removal}
Exact duplicates are removed first. Duplicate keys, meaning the same transaction ID with differing values, are
resolved by keeping the most complete version. IDs are compared after upper-casing and trimming, so that hidden
duplicates such as `` dc2020000001'' are caught. Near-duplicates (same date, product, region, channel and
quantity but different IDs) are only \emph{flagged}. In sales data they are usually genuine repeat orders.

\subsection{Outlier detection and treatment}
Order quantities are strongly right-skewed (skewness \QtySkew; median \QtyMedian units). Each product also has
its own typical order size. A global $z$-score therefore flags too many small-product orders and too few
large-product ones. Four methods were compared in the practice phase: the IQR rule, the $z$-score, the MAD-based
modified $z$-score, and a per-product IQR rule on the log scale. The last of these was the most balanced, and
it is the method used in the pipeline:
\[
  \text{flag}_i = \bigl[\, \log(1+q_i) \notin [\,Q_1^{(p)} - 1.5\,\text{IQR}^{(p)},\; Q_3^{(p)} + 1.5\,\text{IQR}^{(p)}\,] \bigr],
\]
where the quartiles are computed separately for each product $p$. Impossible values (for example, a quantity
below one or a discount above 50\%) are treated as errors: they are blanked and re-derived. Unusual but possible
values are flagged and kept. The real data contained \NOutlierFlags such flags (\OutlierPct\% of orders), all
plausible bulk orders, and no impossible values.

\subsection{Normalisation and encoding}
Scaling is not applied globally. It is applied \emph{inside} each model's pipeline, fitted on the training rows
only, to avoid leakage. Count-like features are log-transformed, $\log(1+x)$, to stabilise variance, and they
are then standardised for scale-sensitive models (linear regression, KNN and SVR). Nominal categories (region,
channel, customer type, category) are one-hot encoded. Month is encoded cyclically as
$(\sin 2\pi m/12, \cos 2\pi m/12)$, so that December and January are neighbours. High-cardinality interactions
(channel $\times$ region, product) are target-encoded with additive smoothing,
\[
  \tilde\mu_k = \frac{n_k \bar y_k + m\,\bar y}{n_k + m}, \qquad m=50,
\]
using training data only. Rare levels are thereby shrunk towards the overall mean.

\subsection{Feature engineering for demand forecasting}
Transactions were aggregated into a \emph{panel} of \PanelSeries products $\times$ months. Each feature for
month $t$ uses only information available at the end of month $t-1$, or information that is planned in advance
(\cref{tab:features}). The same feature function is used for training, testing and future forecasting, so the
rule is enforced in one place. One instructive check compared a correct three-month rolling mean (shifted
before rolling) with a ``leaky'' version that included the current month. The leaky version correlated more
strongly with the target, which is exactly why leakage is tempting and dangerous \citep{kaufman2012}.

\begin{table}[htbp]\centering\small
\caption{Features of the demand-forecasting model (\PanelFeatures in total)}\label{tab:features}
\begin{tabularx}{\linewidth}{@{}lX@{}}
\toprule
\textbf{Group} & \textbf{Features} \\ \midrule
History (lags) & units in months $t-1, t-2, t-3, t-6, t-12$ \\
Rolling statistics & mean of the previous 3, 6 and 12 months; standard deviation of the previous 3 months \\
Momentum & last month's year-on-year growth; 12-month trend ratio \\
Price (planned) & list price of month $t$; year-on-year list-price change \\
Product & category (one-hot); pack weight in grams; gift-pack indicator; months since launch \\
Calendar & month number; cyclical month (sine, cosine); December and February indicators; days in month; time index \\ \bottomrule
\end{tabularx}
\end{table}

\section{Dimensional data store}
The cleaned data are loaded into a star schema in SQLite (\cref{fig:star}). The fact table \texttt{fact\_sales}
holds one row per transaction with its measures. The dimensions describe products, dates (including a
Friday--Saturday weekend flag, as in Bangladesh) and the 220 market segments formed by region, channel and
customer type. Indexes on the foreign keys make typical KPI queries return in well under a second. A
reconciliation check confirms that total net sales in the warehouse equal the total in the cleaned dataset to
the last paisa. \Cref{app:sql} lists example analytical SQL queries.

\begin{figure}[htbp]\centering
\begin{tikzpicture}[every node/.style={font=\scriptsize},
  T/.style={draw=brand, rectangle split, rectangle split parts=2, rounded corners, fill=white, text width=3.4cm, align=left},
  F/.style={T, draw=accent, fill=accent!6, text width=3.8cm}]
\node[F] (f) at (0,0) {\textbf{fact\_sales} (\NRows)\nodepart{second}transaction\_id\\date\_id (FK)\\product\_id (FK)\\segment\_id (FK)\\list/market/unit price\\discount\_pct, quantity\\gross, discount, net sales\\flag\_outlier\_qty};
\node[T] (p) at (-5.2,1.6) {\textbf{dim\_product} (\NProducts)\nodepart{second}product\_id (PK)\\product\_name, category\\pack\_size, launch\_year\\mrp\_bdt};
\node[T] (d) at (5.2,1.6) {\textbf{dim\_date} (2,465)\nodepart{second}date\_id (PK), date\\year, quarter, month\\day\_of\_week, is\_weekend};
\node[T] (s) at (0,-4.1) {\textbf{dim\_segment} (220)\nodepart{second}segment\_id (PK)\\region, channel\\customer\_type};
\draw[thick, brand] (f.west) -- (p.east); \draw[thick, brand] (f.east) -- (d.west); \draw[thick, brand] (f.south) -- (s.north);
\end{tikzpicture}
\caption{Star schema of the analytical data store}\label{fig:star}
\end{figure}
